\cotefile{10}
\scanpage{155}
Soit $G$ un \unsure{groupe|grpe} \deleted{\unsure{quel}} algébrique linéaire quelconque,
$T$ un tore maximal, $N$ son normalisateur et $W=N/T$.
\deleted{\unsure{G}} $N$ opère sur $G$ et $T$ de façon compatible par
\unsure{autom.} intérieurs, et \deleted{\illegible{1}} donc \deleted{\unsure{le}} l'hom canonique
\[ CK(G)\longrightarrow CK(T)\simeq\mathbb Z[\widehat T] \]
est compatible avec les opérations de $N$. $N$ opère
trivialement sur le pr. mb., et $T$ opère trivialement
sur le second, donc en fait $W$ opère sur
$\mathbb Z[\widehat T]$ via les opérations de $W$ sur $\widehat T$. On
a donc un hom
\[ CK(G)\longrightarrow\mathbb Z[\widehat T]^W \tag*{(5)} \]
où $\mathbb Z[\widehat T]^W$ désigne \deleted{le}l'\deleted{groupe}\inserted{anneau} des \unsure{élémts} de
$\mathbb Z[\widehat T]$ invariants sous $W$.

\deleted{On} En général \deleted{\unsure{(à cause de la torsion)}} on n'a
pas là une bijection. Mais c'est \inserted{sans doute} une bijection
s'il n'y a pas de torsion, \deleted{\unsure{(du moins on}} et de
plus l'homomorphisme analogue
\[ C\bigl(K(G)\mathbin{\underset{\mathbb Z}{\otimes}}k\bigr)\longrightarrow k[\widehat T]^W \tag*{(5 bis)} \]
(où $k$ est un corps commutatif) est bijectif
si \deleted{\unsure{G} n'a pas de} \inserted{la car. de $k$ est} première à la torsion de $G$.

% archive sheet number: 154
% paper watermark: ARGENTINE TYPEWRITER (partially visible at the edge)
